\subsection{Phase 5: Prodromal PD Progression Prediction}

\paragraph{Cox Proportional Hazards Model Performance}
The Cox proportional hazards model achieved a concordance index (C-index) of 0.79 (95\% CI: 0.72-0.86) for predicting time to PD diagnosis in the prodromal cohort, significantly outperforming a baseline clinical model using only age, sex, and family history (C-index: 0.64, 95\% CI: 0.56-0.72, $p = 0.003$). Time-dependent area under the ROC curve (AUC) analysis demonstrated robust discrimination at 1-year (AUC: 0.84), 3-year (AUC: 0.82), and 5-year (AUC: 0.78) prediction horizons (Figure 6). Model calibration assessed using Nam-D'Agostino test showed good agreement between predicted and observed conversion probabilities ($\chi^2 = 8.42$, $p = 0.394$).

\paragraph{Key Predictive Features and Hazard Ratios}
Multivariate Cox regression identified several significant predictors of progression to clinical PD (Table 3). \textbf{Baseline motor symptoms} emerged as the strongest predictor: each 5-point increase in MDS-UPDRS Part III score conferred a hazard ratio (HR) of 2.84 (95\% CI: 2.12-3.80, $p < 0.001$). \textbf{Male sex} was associated with increased conversion risk (HR: 1.92, 95\% CI: 1.24-2.97, $p < 0.001$), consistent with PD epidemiology. \textbf{REM sleep behavior disorder (RBD)} substantially elevated risk (HR: 2.31, 95\% CI: 1.58-3.38, $p < 0.001$), while \textbf{hyposmia} showed moderate association (HR: 1.78, 95\% CI: 1.21-2.61, $p = 0.003$).

Among imaging biomarkers, \textbf{reduced DaTscan striatal binding} was strongly predictive: each 0.1-unit decrease in putamen SBR increased risk by 15\% (HR per 0.1 decrease: 1.15, 95\% CI: 1.08-1.23, $p < 0.001$). Structural MRI measures showed that \textbf{substantia nigra volume} was inversely associated with conversion (HR per 100 mm³ decrease: 1.34, 95\% CI: 1.12-1.60, $p = 0.001$), and \textbf{cortical thinning} in entorhinal cortex predicted risk (HR per 0.1 mm decrease: 1.22, 95\% CI: 1.05-1.41, $p = 0.010$).

Genetic factors contributed significant risk: \textbf{\textit{GBA} mutations} (HR: 3.21, 95\% CI: 1.78-5.79, $p < 0.001$), \textbf{\textit{LRRK2} mutations} (HR: 2.47, 95\% CI: 1.32-4.61, $p = 0.005$), and \textbf{\textit{APOE} $\varepsilon$4} (HR: 1.48, 95\% CI: 1.03-2.13, $p = 0.035$) all increased conversion probability. CSF biomarkers showed that \textbf{low $\alpha$-synuclein} (HR per 100 pg/mL decrease: 1.18, 95\% CI: 1.07-1.31, $p = 0.001$) and \textbf{elevated total-tau} (HR per 10 pg/mL increase: 1.14, 95\% CI: 1.04-1.25, $p = 0.006$) predicted progression.

\paragraph{DeepSurv Neural Survival Model}
The DeepSurv neural network architecture, consisting of three fully connected layers with ReLU activation and dropout regularization, achieved a C-index of 0.82 (95\% CI: 0.76-0.88), representing a 3.8\% improvement over the Cox model ($p = 0.041$). This improvement was most pronounced for early conversion prediction (1-year AUC: 0.89 vs. 0.84). Feature importance analysis using gradient-based attribution revealed that DeepSurv learned complex non-linear interactions, particularly between motor symptoms and dopaminergic imaging, and between RBD and $\alpha$-synuclein levels.

\paragraph{Time-Varying Covariate Analysis}
To account for dynamic changes in risk factors, we extended the Cox model to include time-varying covariates for longitudinally measured features. Annual rate of UPDRS-III increase (HR per 1-point/year increase: 1.62, 95\% CI: 1.34-1.96, $p < 0.001$) and rate of DaTscan SBR decline (HR per 0.1-unit/year decrease: 1.41, 95\% CI: 1.19-1.67, $p < 0.001$) substantially improved prediction beyond baseline values alone (C-index increase: 0.79 to 0.84, $p = 0.008$). This suggests that monitoring progression velocity in prodromal individuals enhances risk stratification.

\paragraph{Risk Stratification and Clinical Thresholds}
Based on predicted conversion probabilities, we stratified the prodromal cohort into low-risk (lowest tertile, predicted 3-year conversion probability <15\%), moderate-risk (middle tertile, 15-40\%), and high-risk (upper tertile, >40\%) groups. Observed conversion rates closely matched predictions: low-risk group showed 8.3\% conversion at 3 years, moderate-risk 31.7\%, and high-risk 64.2\% (log-rank $p < 0.001$). These risk categories demonstrated clinical utility by identifying individuals who might benefit from closer monitoring (moderate-risk) or enrollment in disease-modifying therapy trials (high-risk).

For clinical decision-making, we established biomarker thresholds optimized for sensitivity and specificity: UPDRS-III $\geq$ 10 points (sensitivity: 78\%, specificity: 71\%), putamen SBR $\leq$ 1.8 (sensitivity: 82\%, specificity: 65\%), and CSF $\alpha$-synuclein $\leq$ 1,500 pg/mL (sensitivity: 71\%, specificity: 74\%). Combining these three biomarkers achieved 89\% sensitivity and 83\% specificity for identifying individuals who would convert within 3 years, with positive predictive value of 72\% and negative predictive value of 94\%.

\paragraph{Subgroup Performance Analysis}
Model performance varied across demographic and genetic subgroups. The C-index was higher in men (0.81, 95\% CI: 0.73-0.89) than women (0.76, 95\% CI: 0.65-0.87), though not significantly different ($p = 0.318$). Prediction was more accurate in \textit{GBA} mutation carriers (C-index: 0.86, 95\% CI: 0.77-0.95), likely due to more consistent and rapid progression patterns. Performance was similar across age groups (<60 years: 0.78, $\geq$60 years: 0.80, $p = 0.671$), supporting generalizability across the prodromal age spectrum.
